\documentclass{article}
\usepackage[utf8]{inputenc}

\usepackage[a4paper, total={6in, 8in}]{geometry}
\usepackage{amsmath, amssymb, dsfont}
\usepackage{graphicx}

\title{Foundations of Machine Learning -- Homework assignment 1}
\author{CentraleSupélec / Université Paris-Saclay 
\\ Yannick Le Cacheux
\\ yannick.le-cacheux@centralesupelec.fr}
\date{Due date: 15\textsuperscript{th} November 2026}
% \date{Due date: to be announced}

\begin{document}

\maketitle

\paragraph{Instructions.}

This assignment may be subject to small changes regarding the number of points awarded per question.
You are allowed to collaborate with each others to complete the assignment.
However, your solution \emph{must} be your own.
Submission of solutions which are partly not your own, included solutions generated by AI, will result in a 0 for the entire assignment, as per the course rules.

5 bonus points will be granted if your solution is written with \LaTeX~or another suitable typewriting software \emph{and} if you use bold uppercase letters for matrices (\textit{e.g.} $\mathbf{X}$), bold lowercase ones for vectors (\textit{e.g.} $\mathbf{x}$), and regular ones for scalars (\textit{e.g.} $x$).


\section{Fundamental concepts [30 points]}

\paragraph{1.1 Bayes theorem [5 points].} A laboratory has a blood test for a disease which has a \textit{sensitivity} of 0.95, which means that if a tested person has the disease, the test will have a 95\% chance of being positive.
The test also has a \textit{false positive rate} of 1\%, which means that if a tested person does not have the disease, the test will have a 1\% risk of wrongly being positive. If the prevalence of the disease is 0.1\%, i.e. the probability of a random person being infected is 1 out of 1000, what is the probability that a randomly tested person with a positive result is infected?

\paragraph{1.2 Matrix calculus [5 points].} Prove that for a vector $\mathbf{x} \in \mathbb{R}^D$ and a matrix $\mathbf{A} \in \mathbb{R}^{D \times D}$ which does not depend on $\mathbf{x}$,
$$\frac{\partial \mathbf{x}^\top \mathbf{A} \mathbf{x}}{\partial \mathbf{x}} = (\mathbf{A} + \mathbf{A}^\top) \mathbf{x}$$

\paragraph{1.3 Constrained optimization [5 points].} Given functions
$f(x_1, x_2) = x_1^2 + x_2^2$ and $g(x_1, x_2) = x_1 \cdot x_2$, find the minimum of $f(x_1, x_2)$ such that $g(x_1, x_2) = 4$.

\paragraph{1.4 Maximum likelihood estimation.} The normal distribution, also called the Gaussian distribution, is one of the most commonly encountered probability distributions. For given mean $\mu$ and standard deviation $\sigma$, its probability density function is
\begin{equation} \label{eq:normal_distribution}
p(x|\mu, \sigma) = \mathcal{N}(x|\mu, \sigma^2) = \frac{1}{\sqrt{2\pi\sigma^2}} \text{ exp} \left( -\frac{1}{2} \frac{(x-\mu)^2}{\sigma^2}\right)
\end{equation}
We assume that we have $N$ independent and identically distributed (\textit{i.i.d.}) samples $\mathbf{x} = \{x_1, x_2, \dots, x_N\}$ from this distribution. Recall that the \textit{likelihood} of some data $\mathbf{x}$ given parameters $\boldsymbol{\theta}$ is given by $p(\mathbf{x}|\boldsymbol{\theta})$.

\noindent\textbf{1.4.1 [5 points]} Write the log-likelihood of the data, i.e. the logarithm of the likelihood.  You may use a logarithm is base $e$, i.e. the natural logarithm $\text{ln}$.

In general, maximum likelihood is an estimation method for parameters $\boldsymbol{\theta}$ in which we select
$\hat{\boldsymbol{\theta}} = \underset{\boldsymbol{\theta}}{\text{argmax }} p(\mathbf{x}|\boldsymbol{\theta})$.

\noindent\textbf{1.4.2 [5 points]} Setting the gradient of the log-likelihood with respect to $\mu$ to $0$, derive the maximum likelihood estimator of $\mu$.

\noindent\textbf{1.4.3 [5 points]} Using a similar method, derive the maximum likelihood estimator of $\sigma^2$.

\noindent\textbf{1.4.a \textit{Bonus question [+5 points]}}: show that the maximum likelihood estimator of $\sigma^2$ is biased. An estimator $\hat{\boldsymbol{\theta}}$ is unbiased if $\mathbb{E}[\hat{\boldsymbol{\theta}}] = \boldsymbol{\theta}$.


\section{Linear regression [35 points]}

\paragraph{2.1. Parameters of a linear regression.} In linear regression, for a given $D$-dimensional input variable $\mathbf{x} = (x_1 ~x_2~\dots~x_D)^\top \in \mathbb{R}^D$, we estimate a value $y$ with
$\hat{y} = \mathbf{w}^\top \mathbf{x} = \sum_{d=1}^D w_d x_d$. In the least square formulation, we assume that we are given $N$ labeled observations $\{(\mathbf{x}_1, y_1), \dots, (\mathbf{x}_N, y_N)\}$, and we are looking for the parameters $\mathbf{w} \in \mathbb{R}^D$ which minimize
\begin{equation} \label{eq:least_square}
\frac{1}{N} \sum_{n=1}^N (y_n - \mathbf{w}^\top \mathbf{x}_n)^2
\end{equation}
We will assume here that input variables $\mathbf{x}_n$ are 2-dimensional, such that $\mathbf{w} = (w_1~w_2)^\top = \begin{pmatrix} w_1 \\ w_2 \end{pmatrix}$.

\noindent\textbf{2.1.1 [8 points]} Setting the partial derivatives of Equation~(\ref{eq:least_square}) with respect to each parameter $w_1, w_2$ to 0, derive the value of $\mathbf{w}$ given the observations.

Once you are done with question 2.1.1: now, recall from the lecture on linear regression that in general in $D$ dimensions, noting $\mathbf{X} = (\mathbf{x}_1, \dots, \mathbf{x}_N)^\top \in \mathbb{R}^{N \times D}$ our input variables and $\mathbf{y} = (y_1, \dots, y_N) \in \mathbb{R}^N$ our labels, the parameters $\mathbf{w}$ can be obtained with
\begin{equation} \label{eq:least_square_solution}
\mathbf{w} = (\mathbf{X}^\top\mathbf{X})^{-1}\mathbf{X}^\top\mathbf{y}
\end{equation}

\noindent\textbf{2.1.2 [7 points]} Are your results from Question~2.1.1 compatible with Equation~(\ref{eq:least_square_solution})?

\paragraph{2.2 Least square loss.} We will now assume that our input variables have a single dimension. We will further assume that the labels $\{y_1, \dots, y_N\}$ of our observations are given by $y = w \cdot x + \epsilon$, where $\epsilon$ is some random noise sampled from a gaussian distribution with mean zero and (unknown)  standard deviation $\sigma$. We write
$\epsilon \sim \mathcal{N}(0, \sigma^2)$.

\noindent\textbf{2.2.1 [10 points].} Write the log-likelihood of the observations $\{(x_1, y_1), \dots, (x_N,y_N)\}$ as a function of $w$. \textit{Hint}: you may use the fact that for the joint distribution $p(a,b)$ of two random variables $a$ and $b$, we have $p(a,b) = p(a|b)p(b)$.

\noindent\textbf{2.2.2 [10 points].} Show that maximizing the log-likelihood with respect to $w$ is equivalent to minimizing the sum of the squared errors as in Equation~(\ref{eq:least_square}).

\section{Logistic regression [10 points]}

In the gambling world, the ``chance of winning'' a game with two outcomes (gain or loss) are often stated as \textit{odds}, defined as a ratio. For instance, a 7:1 odd means that the probability of winning is 7 times higher than the probability of losing.
As often in machine learning, for values with many different possible orders of magnitudes, it is easier to consider the logarithm of this value.
Assume we want to predict the log of the odd ratio with a linear model, such that
\begin{equation}
\mathbf{w}^\top\mathbf{x} \simeq \text{log} \left( \frac{P(y=1)}{P(y=0)} \right)
\end{equation}
\noindent\textbf{3.1 [5 points].} Prove that predicting the log of the odd ratio with a linear model is equivalent to predicting the probability of winning by applying a sigmoid function to the output of a linear model.

\noindent\textbf{3.2 [5 points].} Prove that training our parameters by maximizing the log-likelihood of a set of observation gives us the logistic loss.
You may want to use the natural logarithm for this exercise.

\section{Clustering [25 points]}

\paragraph{4.1 K-Means [10 points].} In the K-Means algorithm, given $N$ unlabeled points $\{\mathbf{x}_1, \dots, \mathbf{x}_N\}$ and a number of clusters $K$, we apply a greedy algorithm to find $K$ cluster centers $\{\mathbf{c}_1, \dots, \mathbf{c}_K\}$ and a function $a: \mathbb{R}^D \rightarrow \{1, \dots, K\}$ assigning a point $\mathbf{x}$ to one of the $K$ clusters
in order to minimize the inertia
\begin{equation}
\frac{1}{N} \sum_{n=1}^N \lVert \mathbf{x}_n - \mathbf{c}_{a(\mathbf{x}_n)}\rVert_2^2
\end{equation}
For the purpose of this exercise, we can rewrite this objective as
\begin{equation} \label{eq:kmeans_objective}
\frac{1}{N} \sum_{n=1}^N \sum_{k=1}^K a_{n,k} \lVert \mathbf{x}_n - \mathbf{c}_{k}\rVert_2^2
\end{equation}
where $a_{n,k}=1$ if point $\mathbf{x}_n$ is assigned to cluster $k$ and $0$ otherwise.

\noindent Prove that if the cost from Equation~(\ref{eq:kmeans_objective}) is minimized, i.e. we are at the global minimum of the cost, then
\begin{equation}
\mathbf{c}_k = \frac{\sum_{n=1}^N a_{n,k} \mathbf{x}_n}
{\sum_{n=1}^N a_{n,k}}
\end{equation}
That is, the cluster centers $\{\mathbf{c}_1, \dots, \mathbf{c}_K\}$ are the means of the points assigned to the respective clusters.
Note that this result should \textit{not} depend on the actual heuristic (algorithm) we use to minimize to cost function.
\textit{Hint}: you may want to compute the partial derivative of something with respect to something else.

\paragraph{4.2 Hierarchical clustering and Levenshtein distance.} The Levenshtein distance is a string metric measuring the difference between two sequences of characters as the number of character insertions, deletions or substitutions necessary to transform one sequence into the other.

\noindent\textbf{4.2.1 [8 points].} Show that the Levenshtein distance is indeed a distance, using the mathematical definition of the term.

In hierarchical clustering, the Levenshtein distance can be used as an element-wise distance, i.e. as a distance between two points in the feature space. We also need to define a group-wise distance, i.e. a distance between two groups of elements. Given an element-wise distance $d(\mathbf{x}_m,\mathbf{x}_n)$ between two elements $\mathbf{x}_m$ and $\mathbf{x}_n$, the single-linkage criterion $D(\mathcal{C}_i,\mathcal{C}_j)$ between two clusters $\mathcal{C}_i$ and $\mathcal{C}_j$ is defined as
\begin{equation}
D(\mathcal{C}_i, \mathcal{C}_j) =
\underset{\substack{\mathbf{x}_m \in \mathcal{C}_i \\ \mathbf{x}_n \in \mathcal{C}_j}}{\text{min }} d(\mathbf{x}_m,\mathbf{x}_n)
\end{equation}
\noindent\textbf{4.2.2 [7 points].} Is the single-linkage criterion a distance in the general case? Prove the result if it is, and provide a counter-example if it is not.

\end{document}
